\chapter{Contexto e Exemplos: Mostrar a Regra em Ação}
\label{ch:contexto}

\ficcao Chega à biblioteca uma mensagem: «Queria saber se o livro da reserva está pronto e prolongar o empréstimo de outro livro.» A lista de categorias tem «reserva», «renovação» e «outro». Se só podemos guardar uma, a lista deixa por resolver o caso de uma pessoa que pede duas coisas.

\section{Mostrar zero, um ou vários exemplos}

Para este exercício, escolhemos uma regra simples. Se houver um pedido de renovação, assinala «renovação». Se houver só um pedido sobre uma reserva, assinala «reserva». Nos restantes casos, assinala «outro». A prioridade é uma escolha do exercício. Noutro sistema, poderíamos guardar as duas categorias.

Podes dar apenas a regra (\emph{zero-shot}), acrescentar uma mensagem já classificada (\emph{one-shot}) ou mostrar várias (\emph{few-shot}). Uma boa demonstração mostra uma decisão que, de outra forma, seria fácil interpretar de duas maneiras:

\begin{lstlisting}
Mensagem: «A reserva está pronta? Preciso também de
mais uma semana para devolver o livro que levei.»
Categoria: renovação
Justificação: a mensagem também pede para prolongar
um empréstimo, que tem prioridade nesta regra.
\end{lstlisting}

Os exemplos integram o contexto dessa utilização. A resposta pode mudar com a nova informação, mantendo os mesmos parâmetros do modelo. Esta é a ideia de \emph{aprendizagem em contexto} observada nos modelos e tarefas estudados por Brown e colaboradores~\citep{brown2020}. Com as mesmas regras, testa mensagens novas para ver se as demonstrações ajudam no teu caso.

\begin{promptlabbox}{Comparar zero, uma e várias demonstrações}
\textbf{Objetivo e contexto:} comparar três maneiras de apresentar a mesma regra de classificação. Usa mensagens fictícias e escolhe cinco para o teste, incluindo uma com dois pedidos.

\textbf{Procedimento:} classifica as cinco mensagens à mão e guarda essas respostas. Em conversas separadas, fornece a regra e as cinco mensagens. Na primeira, usa só a regra. Na segunda, inclui a demonstração acima. Na terceira, junta duas demonstrações novas, com «reserva» e «outro». Usa o mesmo modelo e a mesma configuração disponível.

\textbf{Resultado esperado e validação:} compara as categorias com as respostas de referência. Regista os erros, incluindo os casos em que as demonstrações introduzem um erro. Cinco mensagens permitem observar comportamentos concretos.

\textbf{Falhas possíveis:} usar uma mensagem do teste como demonstração, mudar a regra entre condições ou contar justificações eloquentes como categorias corretas.
\end{promptlabbox}

\section{O exemplo pode mostrar forma e significado}

Se precisares de uma tabela, escreve o nome das colunas e uma linha completa. Para encaminhar pedidos da biblioteca, «categoria», «excerto» e «informação em falta» permitem guardar a decisão, o trecho que a sustenta e o que ainda falta confirmar. A coluna «excerto» tem de citar palavras que estão na mensagem. A coluna «informação em falta» pode indicar que o número da reserva não foi fornecido. O formato e o significado de cada coluna ficam assim ligados.

Conserva as exceções na regra, mesmo quando todos os exemplos apresentados são simples. Se o modelo propuser mais demonstrações, revê-as com a regra à frente. Um rascunho que dá prioridade a reservas vai ensinar o contrário do que foi pedido. As mensagens usadas para demonstrar ficam separadas das reservadas para avaliar.

\section{Organizar informação que chega de várias fontes}

A \emph{janela de contexto} é o espaço de informação que o modelo pode considerar de uma vez. Um índice curto e títulos claros ajudam o leitor a encontrar a passagem que interessa. Podemos medir se ajudam também o modelo. Quando passares a selecionar excertos automaticamente, o capítulo~\ref{ch:context-engineering} mostra como guardar a origem, a versão e a informação que foi retirada.

\begin{promptlabbox}{Comparar um texto corrido com o mesmo texto organizado}
\textbf{Objetivo e contexto:} perceber se títulos e um índice ajudam a encontrar uma passagem num documento que possas usar.

\textbf{Procedimento:} escolhe uma pergunta factual e sublinha previamente o trecho que a responde. Numa conversa, dá o texto corrido. Noutra, junta títulos e um índice ao mesmo texto. Conserva todas as frases e faz a mesma pergunta, sem dizer em que secção está a resposta.

\textbf{Resultado esperado e validação:} compara a resposta e a localização da evidência. Este ensaio avalia os títulos e o índice em conjunto. Para medir cada contribuição, cria uma condição adicional.

\textbf{Falhas possíveis:} mudar o conteúdo, dar a secção certa numa só condição ou tomar o tom confiante como medida de acerto.
\end{promptlabbox}

Até aqui, demos ao modelo a regra e os exemplos. A seguir, vamos escolher a perspetiva a partir da qual queremos que leia o pedido e distinguir essa escolha dos papéis definidos pela aplicação.
